% -*- coding: utf-8 -*-
\section{信噪比驱动的动态神经门控特征剪枝机制 (DCGU)}
\label{sec:3_2_dcgu_gating}

在第三章第 3.1 节所建立的高速多径时变信道中，车端移动引起的信噪比 $\gamma$ 呈现随机跳变特征。传统的深度信源信道联合编码系统（如经典 Deep JSCC 及变分 NTSCC）通常假定训练信噪比与测试信噪比严格匹配。当测试信道恶化至负信噪比区间时，原先针对高信噪比收敛的网络参数将产生严重的“信噪比失配崩溃”（SNR Mismatch Catastrophe）：解码端将信道强噪声误当作高频语义特征强行逆变换，重构图像产生大面积斑点畸变与伪影。

针对上述难题，本节提出一种基于物理信噪比驱动的超轻量级\textbf{动态通道门控单元（Dynamic Channel Gating Unit, 简称 DCGU）}。该模块赋予发送端在微秒级时间内感知信道波动并自适应切除冗余/脆弱特征通道的自愈能力。

\subsection{门控子网络层级架构与参数轻量化设计}
DCGU 的核心目标是根据瞬时信噪比标量指示因子 $\gamma \in \mathbb{R}$，生成能够精准调制各个潜在特征维度的连续门控权重向量 $\bm{m} \in [0, 1]^K$。

为实现车载边缘端极低的计算开销，门控子网络采用轻量级两层多层感知机（MLP）拓扑。设瓶颈隐藏层维数为 $d_{\text{mid}}$（工程实现中取 $d_{\text{mid}} = 32 \ll K$）。网络前向计算方程表达如下：
\begin{align}
  \bm{u}_1 &= \text{ReLU}(\bm{W}_1 \cdot \gamma + \bm{b}_1) \in \mathbb{R}^{d_{\text{mid}}} \label{eq:dcgu_layer1} \\
  \bm{a}   &= \bm{W}_2 \cdot \bm{u}_1 + \bm{b}_2 \in \mathbb{R}^K \label{eq:dcgu_layer2} \\
  \bm{m}   &= \sigma(\bm{a}) = \left[ \frac{1}{1 + e^{-a_1}}, \frac{1}{1 + e^{-a_2}}, \dots, \frac{1}{1 + e^{-a_K}} \right]^T \label{eq:dcgu_sigmoid}
\end{align}
其中：
\begin{itemize}
  \item $\bm{W}_1 \in \mathbb{R}^{d_{\text{mid}} \times 1}$ 与 $\bm{b}_1 \in \mathbb{R}^{d_{\text{mid}}}$ 分别为第一层线性映射权重与偏置向量；
  \item $\bm{W}_2 \in \mathbb{R}^{K \times d_{\text{mid}}}$ 与 $\bm{b}_2 \in \mathbb{R}^K$ 分别为第二层通道展开投影权重与偏置向量；
  \item $\text{ReLU}(x) = \max(0, x)$ 为非线性修正线性单元，赋予门控网络对信噪比区间的非线性分段拟合能力；
  \item $\sigma(\cdot)$ 为标准 Logistic Sigmoid 激活函数，将门控掩码严格约束至连续概率区间 $(0, 1)$。
\end{itemize}

计算复杂度分析表明，当潜在特征维数 $K = 256$、$d_{\text{mid}} = 32$ 时，整个 DCGU 子网络的浮点参数总量仅为：
\begin{equation}
  N_{\text{DCGU}} = (1 \times 32 + 32) + (32 \times 256 + 256) = 64 + 8448 = 8512 \text{ (参数)}
\end{equation}
以 32 位单精度浮点数存储，所占内存空间不足 $34\text{ KB}$；单次前向推理仅需约 $1.7 \times 10^4$ 次浮点运算（FLOPs），在车载主流嵌入式芯片（如 NVIDIA Orin 或高通 8155）上的执行时延低于 $0.005\text{ ms}$，具备天然的微秒级实时响应优势。

\subsection{软松弛剪枝机制与端到端反向传播梯度流}
若直接采用离散二值硬截断门控（即令 $m_k \in \{0, 1\}$），阶梯函数在零点处的导数为狄拉克脉冲，在非零处导数恒为零，将导致端到端反向传播算法（Backpropagation）无法求解梯度。

为此，本文采用连续可微的软松弛剪枝（Soft Continuous Pruning）协议。特征加权过程形式化为编码特征与门控向量的哈达玛积：
\begin{equation}
  \bm{z}_{\text{pruned}} = \bm{z} \odot \bm{m} \in \mathbb{R}^K
  \label{eq:soft_pruning_formulation}
\end{equation}
展开至标量维度，即：
\begin{equation}
  z_{\text{pruned}, k} = z_k \cdot m_k, \quad \forall k \in \{1, 2, \dots, K\}
\end{equation}

在反向传播训练阶段，设标量损失函数为 $\mathcal{L}$，根据多元微积分链式求导法则，流经 DCGU 的偏导数梯度传播方程解析如下：
\begin{align}
  \frac{\partial \mathcal{L}}{\partial z_k} &= \frac{\partial \mathcal{L}}{\partial z_{\text{pruned}, k}} \cdot m_k \label{eq:grad_to_z} \\
  \frac{\partial \mathcal{L}}{\partial a_k} &= \frac{\partial \mathcal{L}}{\partial z_{\text{pruned}, k}} \cdot z_k \cdot \sigma(a_k)(1 - \sigma(a_k)) \label{eq:grad_to_a} \\
  \frac{\partial \mathcal{L}}{\partial \bm{W}_2} &= \left( \frac{\partial \mathcal{L}}{\partial \bm{a}} \right) \cdot \bm{u}_1^T \label{eq:grad_to_W2} \\
  \frac{\partial \mathcal{L}}{\partial \bm{W}_1} &= \left[ \bm{W}_2^T \left( \frac{\partial \mathcal{L}}{\partial \bm{a}} \right) \odot \mathbb{I}(\bm{u}_1 > 0) \right] \cdot \gamma \label{eq:grad_to_W1}
\end{align}
其中 $\mathbb{I}(\cdot)$ 为指示函数。

公式~(\ref{eq:grad_to_z}) 表明：当信道处于极差状态（如深衰落 $\gamma \to -5\text{ dB}$）时，门控掩码趋向于零（$m_k \to 0$），此时流向特征 $z_k$ 的反向梯度被自动阻断；而在高信噪比时（$m_k \to 1$），梯度完整无损地回传至主编码网络。这种自适应动态阀门设计避免了恶劣信道噪声梯度对主干特征提取器权重造成毁灭性冲击，确保了复杂非凸流形空间的优化平稳性。

\subsection{抗噪物理机制剖析：从信噪比自适应到频域能量重分配}
为什么动态通道门控能从本质上消除传统数字通信的“悬崖效应”？其物理因果机制可通过信源信息熵与物理信道容量匹配定理予以严密阐释：

\begin{enumerate}
  \item \textbf{高信噪比工作区（$\gamma > 12\text{ dB}$，信道容量充裕）}：
  门控网络输出 $\bm{m} \approx [1, 1, \dots, 1]^T$。全部 $K$ 个潜在特征通道全部激活并接入信道。编码器不仅传输大尺度宏观语义结构（低频分量），同时完整下发高频纹理与细节轮廓信息，重构图像 PSNR 达到 $36\text{ dB}$ 以上的高保真水平；
  \item \textbf{低信噪比工作区（$-5\text{ dB} \le \gamma < 5\text{ dB}$，信道受限深衰落）}：
  物理信道香农容量发生剧烈萎缩。若强行传输全部 $K$ 个符号，有限的总发射功率 $P$ 将被均分至所有维度，导致每个符号维度的平均接收能量低于噪声基底。此时，DCGU 主动将抗噪性能最脆弱的高频纹理通道权重压制至零（$m_k \approx 0$），有效活跃通道数压缩至 $K_{\text{eff}} \approx 0.35 K$；
  \item \textbf{发射能量的自适应集中放大效应}：
  结合第 3.1 节公式~(\ref{eq:power_normalization_rigorous}) 的功率归一化算子，当大量次要通道被置零后，发射功率集中分配至少数主干语义通道：
  \begin{equation}
    \tilde{z}_{\text{pruned}, k} = \sqrt{K P} \frac{z_{\text{pruned}, k}}{\sqrt{\sum_{j=1}^{K_{\text{eff}}} z_{\text{pruned}, j}^2}}
  \end{equation}
  由于分母中的非零模长显著减小，核心语义特征的符号幅度获得等效增益放大，确保了核心语义信息能够硬性穿透极端恶劣信道，重构图像质量实现优雅平滑降级（Graceful Degradation）。
\end{enumerate}

表~\ref{tab:dcgu_complexity_spec} 详细列出了 DCGU 模块与主流基准模块在参数规模、计算延时与显存消耗层面的实测对比，证实了其在硬件层面的极致轻量性。

\begin{table}[htbp]
  \centering
  \small
  \caption{动态通道门控单元 (DCGU) 与主流注意力/特征提取模块计算复杂度横向对比表}
  \label{tab:dcgu_complexity_spec}
  \resizebox{\textwidth}{!}{%
  \begin{tabular}{l c c c c c}
    \toprule
    \textbf{特征处理模块架构} & \textbf{可学习参数量} & \textbf{单次推理 FLOPs} & \textbf{边缘推理耗时 (ms)} & \textbf{信道自适应能力} & \textbf{信噪比感知机制} \\
    \midrule
    标准多头自注意力 (MHSA)   & $6.29 \times 10^5$ & $4.18 \times 10^7$ & $3.450$ & 无 (固定拓扑) & 不支持 \\
    Squeeze-and-Excitation (SE) & $3.28 \times 10^4$ & $6.55 \times 10^4$ & $0.042$ & 仅依赖图像信源 & 不支持信道输入 \\
    变分超先验熵网络 (NTSCC)   & $1.45 \times 10^6$ & $9.82 \times 10^7$ & $8.920$ & 隐式连续流形 & 需繁杂迭代反演 \\
    \midrule
    \textbf{本文所提 DCGU 门控} & $\mathbf{8.51 \times 10^3}$ & $\mathbf{1.70 \times 10^4}$ & $\mathbf{0.005}$ & \textbf{微秒级全动态} & \textbf{端到端物理驱动} \\
    \bottomrule
  \end{tabular}%
  }
  \vspace{0.15cm}\par
  \noindent\footnotesize{
    \textbf{注}：基准硬件平台统一为 NVIDIA Jetson AGX Orin 嵌入式平台，批处理大小 $B=1$，测试输入特征维数 $K=256$。
  }
\end{table}
